\documentclass[12pt]{article}
\usepackage[utf8]{inputenc}
\usepackage[T1]{fontenc}
\usepackage{hyperref}
\title{Hurrian Language Resources for K-12 Curriculum Development}
\author{Compiled by Claude Sonnet 5.5 (claude-sonnet-5-5) for Andr\'as Kornai}
\date{}
\begin{document}
\maketitle

\section{Introduction}
Hurrian (ISO 639-3 xhu; Glottocode hurr1240; Hurro-Urartian family) was spoken in the ancient Near East and is attested from roughly the third to the late second millennium BCE. Its closest known relative is Urartian. It has no native speakers today. The best-known text is the Mitanni Letter, a roughly 500-line letter from King Tushratta to Amenhotep III, found at Amarna in 1887. Other Hurrian material comes from Hattusa, Mari, Ugarit and other sites. I found no school programme or revitalization community.

\section{Primary Educational Resources}
I found no primers, textbooks or school materials for Hurrian aimed at pupils. The resources below are scholarly or popular overviews that a teacher could adapt.

\section{Linguistic and Reference Materials}
\begin{itemize}
\item Wikipedia, ``Hurrian language'': \url{https://en.wikipedia.org/wiki/Hurrian_language}
\item Wikipedia, ``Mitanni Letter'' (the main long Hurrian text): \url{https://en.wikipedia.org/wiki/Mitanni_Letter}
\item Giorgieri, ``Schizzo grammaticale della lingua Hurrica'' (grammar sketch, abstract at Urkesh.org): \url{https://urkesh.org/abstracts/913Giorgieri2000a.htm}
\item Further Giorgieri abstract at Urkesh.org: \url{https://urkesh.org/abstracts/923Giorgieri2000a.htm}
\item Mark Ronan, ``Hurrian and the Hurrians'' (popular overview): \url{https://www.markronan.com/languages/hurrian/}
\end{itemize}
Search summaries also name Ilse Wegner's \emph{Einf\"uhrung in die Hurritische Sprache} (Harrassowitz, 2000) as the most recent comprehensive grammar. I saw no URL for it, so it is not in the JSON list. I found no dedicated Hurrian dictionary in the results.

\section{Religious and Community Resources}
Much Hurrian text from Hattusa is ritual and religious in character. I found no living religious or community organisation for Hurrian.

\section{Digital and Audio Resources}
I found no audio, video or software resources for Hurrian. Any pronunciation would be scholarly reconstruction.

\section{Curriculum Development Context}
Material is limited. A K-12 unit could use the Mitanni Letter and the history of the Hurrians and Mitanni as content, mostly in English. A few reconstructed Hurrian words could be added with a clear caveat. Ancient-history and linguistics electives are the most plausible places for it.

\end{document}